% ============================================================
%  CAPÍTULO 6 — IMPLEMENTACIÓN
% ============================================================
\chapter{Implementación}
\label{ch:implementacion}

\section{Organización del código}
\label{sec:org-codigo}

El proyecto se organiza en dos repositorios:

\begin{description}
  \item[\texttt{uniback}] Monorepo Python que contiene el framework backend:
    \texttt{src/uniback/} (núcleo), \texttt{tests/} y ficheros de configuración Docker.
  \item[\texttt{gui-components}] Monorepo Angular con la librería (\texttt{projects/ngt-gui/})
    y la aplicación de demostración (\texttt{projects/ngt-gui/gui/}).
\end{description}

\section{Backend — uniback}
\label{sec:impl-backend}

\subsection{Estructura de módulos}
\label{subsec:modulos-backend}

\begin{verbatim}
src/uniback/
  api/
    routers/
      gui.py          # Endpoints de pantallas dinámicas
    crud_factory.py   # Fábrica de routers CRUD genéricos
  core/
    schema/
      generator.py    # Punto de entrada de la tubería
      plugins/        # Plugins de anotación
        base.py
        label.py
        foreign_key.py
        validator.py
        widget.py
        readonly.py
    models/           # Modelos SQLAlchemy de la aplicación
    database.py       # Configuración de sesión SQLAlchemy
  main.py             # Punto de entrada FastAPI
\end{verbatim}

\subsection{Generación de esquemas}
\label{subsec:generacion-esquemas}

El punto de entrada de la tubería es la función \texttt{generate\_schema(model)}, que
aplica los plugins registrados en orden de prioridad:

\begin{lstlisting}[language=Python, caption={Función principal de generación de esquemas.}]
def generate_schema(
    model: type[DeclarativeBase],
    plugins: list[SchemaPlugin] | None = None,
) -> dict:
    # 1. Pydantic schema base
    pydantic_model = sqlalchemy_to_pydantic(model)
    schema = pydantic_model.model_json_schema()

    # 2. Aplicar plugins en orden de prioridad
    effective_plugins = sorted(
        plugins or _DEFAULT_PLUGINS, key=lambda p: p.priority
    )
    for col_name, col_info in inspect(model).mapper.column_attrs.items():
        for plugin in effective_plugins:
            schema["properties"][col_name] = plugin.apply(
                schema["properties"].get(col_name, {}),
                model, col_name, col_info,
            )

    return schema
\end{lstlisting}

\subsection{Fábrica de routers CRUD}
\label{subsec:crud-factory}

La función \texttt{make\_simple\_rest\_crud(model, prefix)} genera un router FastAPI
completo con los endpoints estándar (\texttt{GET /}, \texttt{GET /:id},
\texttt{POST /}, \texttt{PUT /:id}, \texttt{DELETE /:id}) y el endpoint de
operaciones en bloque (\texttt{POST /bulk}).

\begin{lstlisting}[language=Python, caption={Registro de una entidad con la fábrica CRUD.}]
# En main.py o en el módulo de la entidad
router = make_simple_rest_crud(
    model=Identity,
    prefix="/api/identities",
    schema_overrides={"name": {"x-label": "Nombre completo"}},
)
app.include_router(router)
\end{lstlisting}

\subsection{Endpoint de operaciones en bloque}
\label{subsec:impl-bulk}

La implementación del endpoint \texttt{/bulk} distingue tres tipos de operación según
los campos presentes en cada ítem del array de entrada:

\begin{lstlisting}[language=Python, caption={Lógica de clasificación de operaciones bulk.}]
for idx, item in enumerate(items):
    op = item.get("_op")
    item_id = item.get("id")
    data = {k: v for k, v in item.items()
            if k not in {"id", "_op"}}
    try:
        if op == "delete":
            _bulk_delete(db, model, item_id)
        elif item_id is not None:
            _bulk_update(db, model, item_id, data, schema)
        else:
            _bulk_create(db, model, data, schema)
        results.append({"ok": True, "id": item_id})
    except Exception as exc:
        results.append({"ok": False, "errors": _format_error(exc)})
        if on_error == "abort":
            raise HTTPException(status_code=422, ...)
\end{lstlisting}

El modo \textbf{partial} utiliza savepoints de SQLAlchemy
(\texttt{db.begin\_nested()}) para aislar cada operación dentro de la transacción
global:

\begin{lstlisting}[language=Python, caption={Uso de savepoints para el modo partial.}]
sp = db.begin_nested()
try:
    _execute_operation(db, ...)
    sp.commit()
    results.append({"ok": True})
except Exception as exc:
    sp.rollback()
    results.append({"ok": False, "errors": _format_error(exc)})
\end{lstlisting}

\subsection{Endpoint de pantallas dinámicas}
\label{subsec:impl-screens}

El endpoint \texttt{GET /screens/<screen\_name>} construye la definición de pantalla
a partir del registro de modelos y el tipo de pantalla solicitado:

\begin{lstlisting}[language=Python, caption={Generación de una pantalla de tipo editable\_table.}]
def _synthetic_editable_table(model, schema, endpoint) -> dict:
    columns = [
        _build_column_def(name, prop, schema)
        for name, prop in schema["properties"].items()
        if name not in _INTERNAL_FIELDS
    ]
    return {
        "type": "editable_table",
        "title": model.__tablename__.replace("_", " ").title(),
        "fieldId": "id",
        "endpoint": endpoint,
        "bulkEndpoint": f"{endpoint}/bulk",
        "columns": columns,
        "actions": ["create", "edit", "delete"],
        "pagination": {"defaultPageSize": 100},
        "synthetic": True,
    }
\end{lstlisting}

\section{Frontend — ngt-gui}
\label{sec:impl-frontend}

\subsection{Mapeo JSON Schema → Formly}
\label{subsec:schema-to-formly}

El módulo \texttt{schema-to-formly.ts} convierte una propiedad JSON~Schema en un campo
Formly. La función principal inspecciona el tipo base (\texttt{string}, \texttt{integer},
\texttt{boolean}) y las extensiones \texttt{x-*} para seleccionar el tipo de widget:

\begin{lstlisting}[language=TypeScript, caption={Fragmento del mapeo JSON Schema $\to$ Formly.}]
export function propertyToFormlyField(
  name: string,
  prop: JsonSchemaProperty,
): FormlyFieldConfig {
  const widget = prop['x-widget'] ?? inferWidget(prop);
  return {
    key: name,
    type: widget,
    props: {
      label: prop['x-label'] ?? name,
      required: prop['x-required'] ?? false,
      readonly: prop['x-readonly'] ?? false,
    },
    validators: buildValidators(prop),
  };
}
\end{lstlisting}

\subsection{Mapeo JSON Schema → AG Grid}
\label{subsec:schema-to-aggrid}

El módulo \texttt{schema-to-aggrid.ts} convierte la lista de columnas de una definición
de pantalla en definiciones de columna AG~Grid (\texttt{ColDef}). Para cada columna
editable se genera un editor en función del tipo declarado:

\begin{lstlisting}[language=TypeScript, caption={Mapeo de tipo a editor AG~Grid.}]
function resolveEditor(col: ScreenEditableColumn): string | undefined {
  if (col.options || col.optionsEndpoint) return 'agSelectCellEditor';
  switch (col.type) {
    case 'integer':
    case 'number':  return 'agNumberCellEditor';
    case 'boolean': return 'agCheckboxCellEditor';
    default:        return 'agTextCellEditor';
  }
}
\end{lstlisting}

\subsection{BulkEditGridComponent}
\label{subsec:bulk-edit-grid}

El componente \texttt{BulkEditGridComponent} implementa la cuadrícula de edición masiva.
Sus responsabilidades principales son:

\begin{enumerate}
  \item \textbf{Carga inicial:} petición \texttt{GET /<entityPath>/} para obtener los
    registros existentes.
  \item \textbf{Gestión de estado de fila:} cada fila mantiene un campo interno
    \texttt{\_rowState} con valores \texttt{unchanged}, \texttt{new},
    \texttt{modified} o \texttt{deleted}.
  \item \textbf{Validación local:} la función \texttt{validateCell()} aplica las
    restricciones del esquema antes de permitir el envío.
  \item \textbf{Envío en bloque:} construcción del payload para \texttt{POST /bulk}
    y aplicación de los resultados del backend (actualización de IDs, marcado de
    errores por celda).
\end{enumerate}

Las filas con estado distinto de \texttt{unchanged} reciben clases CSS específicas
(\texttt{row-state-new}, \texttt{row-state-modified}, \texttt{row-state-deleted}) para
su diferenciación visual.

\subsection{DynamicPageComponent}
\label{subsec:dynamic-page}

\texttt{DynamicPageComponent} es el componente raíz que enruta las vistas dinámicas.
Tras recibir la definición de pantalla del resolver, selecciona el componente hijo
adecuado mediante \texttt{*ngIf} sobre \texttt{screenType}:

\begin{lstlisting}[language=html, caption={Selección de vista en dynamic-page.component.html (fragmento).}]
<app-dynamic-browse
  *ngIf="screenType === 'browser'"
  [BrowserType]="browserType"
  [HiddenBulkInsert]="false"
  (BulkInsertEvent)="onBulkInsert()">
</app-dynamic-browse>

<ngt-bulk-edit-grid
  *ngIf="screenType === 'editable_table'"
  [entityPath]="bulkEntityPath"
  [screenColumns]="bulkColumns">
</ngt-bulk-edit-grid>
\end{lstlisting}

\subsection{Formly Editor}
\label{subsec:formly-editor}

\todo[inline]{Completar esta sección cuando el Formly Editor esté terminado (Iteración 4).}

El editor visual de esquemas permite al usuario modificar las propiedades de los campos
de un formulario de forma interactiva. Se compone de tres paneles:

\begin{description}
  \item[Árbol de campos (izquierda):] lista jerárquica de los campos del esquema con
    indicadores de tipo y estado.
  \item[Panel de propiedades (centro):] formulario para editar las propiedades del campo
    seleccionado (\texttt{x-label}, \texttt{x-widget}, \texttt{x-required},
    validadores, \ldots).
  \item[Vista previa en vivo (derecha):] formulario Formly generado en tiempo real a
    partir del esquema editado.
\end{description}

Los cambios se persisten en el backend mediante una llamada \texttt{PATCH} al endpoint
de pantalla, que actualiza los \textit{schema overrides} asociados a la entidad.

\section{Integración Docker}
\label{sec:docker}

El sistema se despliega mediante un fichero \texttt{docker-compose.yml} que define tres
servicios:

\begin{description}
  \item[\texttt{db}:] PostgreSQL~16 con volumen persistente y variables de entorno para
    credenciales.
  \item[\texttt{backend}:] imagen Python construida desde el \texttt{Dockerfile} del
    repositorio \textit{uniback}; monta el código fuente en modo desarrollo para
    \textit{hot-reload}.
  \item[\texttt{frontend}:] imagen Node.js que ejecuta \texttt{ng serve}; proxifica las
    peticiones \texttt{/api/*} al servicio \texttt{backend}.
\end{description}

\begin{lstlisting}[language=yaml, caption={Fragmento del docker-compose.yml de desarrollo.}]
services:
  backend:
    build: ./Back/uniback
    environment:
      DATABASE_URL: postgresql+asyncpg://user:pass@db:5432/uniback
      SECRET_KEY: ${SECRET_KEY}
    volumes:
      - ./Back/uniback/src:/app/src
    depends_on: [db]

  frontend:
    build: ./Front/gui-components
    environment:
      NG_PROXY_TARGET: http://backend:8000
    ports:
      - "4200:4200"
    depends_on: [backend]
\end{lstlisting}
